% !TEX encoding = UTF-8
\chapter{Hướng dẫn vận hành và Cơ chế hoạt động của các dịch vụ}
\label{ch:appendix}

Phụ lục này mô tả chi tiết cơ chế hoạt động, các đoạn mã cấu hình minh họa và kết quả chạy lệnh thực tế của từng dịch vụ nghiệp vụ và bảo mật trong hệ thống Job7189, đồng thời cung cấp hướng dẫn vận hành chi tiết cho quản trị viên hệ thống.

\section{Cơ chế hoạt động của các dịch vụ}

Hệ thống được tổ chức thành các phân vùng mạng độc lập ở cấp độ không gian tên (namespace) để thực thi ranh giới bảo mật nghiêm ngặt. Dưới đây là cơ chế hoạt động, cấu hình minh họa và kết quả vận hành của từng dịch vụ.

\subsection{Cổng giao tiếp Kong (Kong Gateway)}
Cổng giao tiếp Kong đóng vai trò là điểm kiểm soát duy nhất cho luồng dữ liệu đi vào hệ thống từ mạng ngoài (luồng Bắc-Nam).
\begin{itemize}
    \item \textbf{Cơ chế hoạt động}: Cổng giao tiếp chạy ở chế độ không lưu trữ trạng thái cơ sở dữ liệu (declarative/DB-less mode) thông qua tệp cấu hình nguồn khai báo \texttt{kong.yml}. Cổng thực hiện đối chiếu mã thông báo bảo mật JWT từ tiêu đề \texttt{Authorization} gửi kèm trong yêu cầu HTTP của người dùng bằng khóa công khai tương ứng. Hệ thống sử dụng một trình cắm tiền chức năng (pre-function plugin) viết bằng ngôn ngữ Lua để thực hiện giải mã mã thông báo JWT và gửi thông tin yêu cầu đến bộ phân quyền OPA nhằm đưa ra quyết định cho phép hoặc từ chối yêu cầu truy cập trước khi chuyển tiếp đến các vi dịch vụ phía sau.
    \item \textbf{Cấu hình định tuyến và trình cắm Lua}: Đoạn cấu hình khai báo dịch vụ nghiệp vụ và trình cắm tiền chức năng được mô tả như sau:
\begin{minted}[fontsize=\footnotesize]{yaml}
# Trích từ cấu hình nguồn infras/kong/kong.yml
services:
  - name: identity-service
    url: http://identity-service.job7189-apps.svc.cluster.local:80
    routes:
      - name: identity-profiles
        paths: ["/api/recruiters/profile"]
        strip_path: false
        plugins:
          - name: jwt
            config:
              claims_to_verify: [exp]

plugins:
  - name: pre-function
    config:
      access:
        - |
          local cjson   = require("cjson.safe")
          local http    = require("resty.http")
          if kong.request.get_method() == "OPTIONS" then return end
          
          local OPA_URL     = "https://opa.security.svc.cluster.local:8181/v1/data/zta/authz/allow"
          local req_method  = kong.request.get_method()
          local req_path    = ngx.var.uri or "/"
          local auth        = kong.request.get_header("authorization") or ""
          local token       = auth:match("^[Bb]earer%s+(.+)$") or ""
          
          local function b64url_decode(s)
            s = s:gsub("-", "+"):gsub("_", "/")
            local pad = 4 - (#s % 4)
            if pad < 4 then s = s .. string.rep("=", pad) end
            return ngx.decode_base64(s)
          end
          
          local claims = {}
          if token ~= "" then
            local _, payload_b64, _ = token:match("([^.]+)%.([^.]+)%.([^.]+)")
            if payload_b64 then
              local raw = b64url_decode(payload_b64)
              if raw then claims = cjson.decode(raw) or {} end
            end
          end
          
          local segments = {}
          for seg in req_path:gmatch("[^/]+") do table.insert(segments, seg) end
          
          local opa_input = {
            input = {
              method = req_method,
              path = req_path,
              path_segments = segments,
              jwt = claims,
            }
          }
          
          local httpc = http.new()
          httpc:set_timeout(1000)
          local res, err = httpc:request_uri(OPA_URL, {
            method = "POST",
            body = cjson.encode(opa_input),
            headers = { ["Content-Type"] = "application/json" },
            ssl_verify = false,
          })
          
          if not res or res.status ~= 200 then return end
          local decoded = cjson.decode(res.body or "") or {}
          if decoded.result ~= true then
            return kong.response.exit(403, {
              message = "forbidden",
              reason = "OPA denied user-authz",
              user = claims.preferred_username or "anonymous"
            })
          end
\end{minted}
    \item \textbf{Kết quả chạy lệnh minh họa}: Khi gửi yêu cầu không có mã thông báo hoặc khi mã thông báo bị OPA từ chối:
\begin{minted}[fontsize=\footnotesize]{bash}
# Trường hợp không gửi mã thông báo JWT (Kong chặn trực tiếp tại route):
$ curl -i -H "Host: api.job7189.local" http://100.108.231.127:30000/api/recruiters/profile
HTTP/1.1 401 Unauthorized
Content-Type: application/json
Connection: keep-alive
Server: kong/3.6

{"message":"Unauthorized"}

# Trường hợp gửi mã thông báo hợp lệ nhưng bị chính sách OPA từ chối truy cập:
$ curl -i -H "Host: api.job7189.local" -H "Authorization: Bearer <valid_jwt>" http://100.108.231.127:30000/api/admin/users
HTTP/1.1 403 Forbidden
Content-Type: application/json
Connection: keep-alive
Server: kong/3.6

{"message":"forbidden","reason":"OPA denied user-authz","user":"recruiter1"}
\end{minted}
\end{itemize}

\subsection{Dịch vụ xác thực Keycloak}
Keycloak đóng vai trò là trung tâm định danh người dùng cuối, chịu trách nhiệm quản lý thông tin tài khoản của nhà tuyển dụng và ứng viên.
\begin{itemize}
    \item \textbf{Cơ chế hoạt động}: Keycloak thực hiện xác thực người dùng dựa trên cấu hình phân hệ (realm) của dự án. Sau khi xác thực thành công, Keycloak ký số mã thông báo truy cập dạng JWT bằng khóa riêng tư của mình. Khóa công khai dùng để xác thực chữ ký được công khai tại đường dẫn JWKS để cổng giao tiếp Kong và các vi dịch vụ tải về đối chiếu cục bộ.
    \item \textbf{Kết quả chạy lệnh lấy mã thông báo truy cập}: Quản trị viên hoặc ứng dụng khách có thể lấy mã thông báo truy cập bằng giao thức OpenID Connect qua luồng mật khẩu:
\begin{minted}[fontsize=\footnotesize]{bash}
# Gửi yêu cầu lấy mã thông báo truy cập JWT từ Keycloak
$ curl -s -X POST "http://100.108.231.127:30000/realms/job7189/protocol/openid-connect/token" \
  -H "Content-Type: application/x-www-form-urlencoded" \
  -d "grant_type=password" \
  -d "client_id=recruiter-app" \
  -d "username=recruiter1" \
  -d "password=recruiter1" | jq

# Kết quả phản hồi chứa mã thông báo truy cập:
{
  "access_token": "eyJhbGciOiJSUzI1NiIsInR5cCIgOiAiSldUIiwia2lkIiA6ICJ4Nk...",
  "expires_in": 300,
  "refresh_expires_in": 1800,
  "refresh_token": "eyJhbGciOiJSUzI1NiIsInR5c...",
  "token_type": "Bearer",
  "not-before-policy": 0,
  "session_state": "a3d7b489-0112-4211-9a74-0fc03bf6ab18",
  "scope": "openid profile email"
}
\end{minted}
    \item \textbf{Kết quả chạy lệnh truy vấn thông tin khóa công khai (JWKS)}: Các dịch vụ có thể kiểm tra danh sách khóa công khai đang hoạt động tại Keycloak:
\begin{minted}[fontsize=\footnotesize]{bash}
$ curl -s http://keycloak.security.svc.cluster.local:8080/realms/job7189/protocol/openid-connect/certs | jq
{
  "keys": [
    {
      "kid": "4x60x9miRRead...",
      "kty": "RSA",
      "alg": "RS256",
      "use": "sig",
      "n": "u1L7Z2m...",
      "e": "AQAB"
    }
  ]
}
\end{minted}
\end{itemize}

\subsection{Hệ thống định danh vi dịch vụ SPIRE/SPIFFE}
SPIRE cung cấp cơ chế định danh tự động và bảo mật cho từng tiến trình chạy trong container (workload) thông qua chứng thư định danh SVID dạng X.509 hoặc JWT.
\begin{itemize}
    \item \textbf{Cơ chế hoạt động}: Tác nhân SPIRE (SPIRE Agent) lắng nghe trên một socket Unix cục bộ trên mỗi node. Khi Pod khởi chạy, SPIRE Agent truy vấn thông tin định danh của Pod từ hệ thống Kubernetes (chẳng hạn không gian tên, tài khoản dịch vụ, nhãn dán của Pod) để thực hiện xác thực thực thể, sau đó cấp phát chứng thư định danh SVID tương ứng cho Pod thông qua Workload API được gắn kết (mount) vào Pod bằng trình điều khiển CSI (\texttt{csi.spiffe.io}).
    \item \textbf{Cấu hình ClusterSPIFFEID}: Đoạn mã dưới đây mô tả chính sách đăng ký định danh tự động cho các dịch vụ nghiệp vụ:
\begin{minted}[fontsize=\footnotesize]{yaml}
# Trích từ cấu hình nguồn infras/k8s-yaml/spire/cluster-spiffe-ids.yaml
apiVersion: spire.spiffe.io/v1alpha1
kind: ClusterSPIFFEID
metadata:
  name: zta-default-workload-identity
spec:
  spiffeIDTemplate: "spiffe://zta.job7189/ns/{{ .PodMeta.Namespace }}/sa/{{ .PodSpec.ServiceAccountName }}"
  podSelector:
    matchExpressions:
      - key: cilium.zta/role
        operator: Exists
  namespaceSelector:
    matchExpressions:
      - key: kubernetes.io/metadata.name
        operator: In
        values:
          - job7189-apps
          - gateway
          - vault
          - security
          - monitoring
          - data
          - management
  dnsNameTemplates:
    - "{{ index .PodMeta.Labels \"app\" }}.{{ .PodMeta.Namespace }}.svc.cluster.local"
\end{minted}
    \item \textbf{Kết quả chạy lệnh minh họa}: Kiểm tra danh sách các thực thể vi dịch vụ được đăng ký trên Máy chủ SPIRE (SPIRE Server):
\begin{minted}[fontsize=\footnotesize]{bash}
$ kubectl exec -n spire spire-server-0 -c spire-server -- /opt/spire/bin/spire-server entry show
Found 50 entry(s)
Entry ID      : 019eb4d9-4dd7-728e-af71-09cbcf7fb555
SPIFFE ID     : spiffe://zta.job7189/ns/job7189-apps/sa/identity-service
Parent ID     : spiffe://zta.job7189/spire/agent/k8s_psat/7189srv05
Regis. delay  : 0s
Selectors     : k8s:ns:job7189-apps
Selectors     : k8s:sa:identity-service
Selectors     : k8s:pod-label:app:identity-service
\end{minted}
\end{itemize}

\subsection{Hệ thống quản lý khóa động Vault}
Vault quản lý toàn bộ các thông tin mật tĩnh và cấp phát tài khoản truy cập cơ sở dữ liệu động có thời hạn ngắn.
\begin{itemize}
    \item \textbf{Cơ chế hoạt động}: Thay vì sử dụng phương thức xác thực Kubernetes truyền thống, hệ thống đã cấu hình phương thức xác thực \texttt{auth/jwt}. Khi vi dịch vụ nghiệp vụ như \texttt{identity-service} khởi chạy, tiến trình phụ \texttt{spiffe-helper} trong Pod sẽ lấy mã thông báo định danh JWT-SVID từ SPIRE Agent và ghi vào một tệp cục bộ trên bộ nhớ đệm. Tiến trình \texttt{vault-agent} sau đó sử dụng mã thông báo này để gửi yêu cầu đăng nhập đến Vault thông qua đường dẫn \texttt{auth/jwt/login} dưới vai trò (\texttt{role}) tương ứng. Khi xác thực thành công, Vault Agent lấy và lưu trữ thông tin tài khoản truy cập động MySQL được sinh ra tự động với thời hạn sử dụng tối đa là 3600 giây (1 giờ).
    \item \textbf{Cấu hình chính sách Vault}: Đoạn mã cấu hình quyền hạn cho phép vi dịch vụ nghiệp vụ đọc thông tin tài khoản động:
\begin{minted}[fontsize=\footnotesize]{text}
# Vault Policy: identity-service-policy
path "database/creds/identity-service" {
  capabilities = ["read"]
}
path "secret/data/production/identity-service" {
  capabilities = ["read"]
}
\end{minted}
    \item \textbf{Kết quả chạy lệnh minh họa}: Mô phỏng câu lệnh đăng nhập và đọc thông tin tài khoản động MySQL được thực hiện tự động:
\begin{minted}[fontsize=\footnotesize]{bash}
# Đăng nhập vào Vault bằng mã thông báo định danh JWT-SVID
$ vault write auth/jwt/login jwt=@/tmp/jwt_svid.token role=identity-service
Key                  Value
---                  -----
token                hvs.CAESIKP_w8qG4vA4Ld_XmP...
token_accessor       Z9O3sV3hY45Y1e...
token_duration       20m
token_renewable      true
token_policies       ["default" "identity-service"]
identity_policies    []
policies             ["default" "identity-service"]

# Đọc tài khoản cơ sở dữ liệu động từ Vault
$ vault read database/creds/identity-service
Key                Value
---                -----
lease_id           database/creds/identity-service/hvs.Pap9ErlQGUCRSyhrZtDue5iq.lease-abc123xyz
lease_duration     3600
lease_renewable    true
password           A1b2C3d4_E5f6G7
username           v-token-identity-service-abcdef1234
\end{minted}
\end{itemize}

\subsection{Bộ điều khiển quyết định chính sách PDP (Policy Decision Point)}
Bộ điều khiển PDP đóng vai trò là bộ não thích ứng, liên tục điều chỉnh nhãn phân hạng mức độ tin cậy dựa trên lỗ hổng bảo mật được phát hiện thực tế.
\begin{itemize}
    \item \textbf{Cơ chế hoạt động}: Bộ điều khiển PDP định kỳ quét các báo cáo lỗ hổng an ninh từ Trivy Operator để tính toán điểm tin cậy và gán nhãn phân hạng mức độ tin cậy (\texttt{high}, \texttt{medium}, \texttt{low}) trực tiếp lên Pod. Nhãn dán này sau đó được Cilium sử dụng làm thuộc tính mạng để điều chỉnh quyền truy cập mạng.
    \item \textbf{Kết quả chạy lệnh minh họa}: Xem danh sách các nhãn tin cậy được PDP gán động trên các Pod ứng dụng trong cụm:
\begin{minted}[fontsize=\footnotesize]{bash}
$ kubectl get pods -n job7189-apps -o jsonpath='{range .items[*]}{.metadata.name}: {.metadata.labels.zta\.job7189/score-bucket}{"\n"}{end}'
candidate-service-7ffc7b57b7-22shz: high
candidate-service-7ffc7b57b7-v4j7w: high
candidate-service-redis-5b484c7f88-mcj7q: high
communication-service-6dc7d6dd46-f8mv4: low
communication-service-redis-fc6645965-hhrkn: high
hiring-service-85c57cc589-m9zs6: high
hiring-service-redis-6585c68cb9-87q8b: high
identity-service-657f64c7d5-ks7x4: high
identity-service-redis-569fb4c444-mgngz: high
job-service-7755d94d47-gbftq: high
job-service-redis-5b987-qmjvw: high
storage-service-7c5db66b7-nsjk4: low
storage-service-redis-78cb79b765-6zkfm: high
workspace-service-6c68c569d9-4frrg: high
workspace-service-redis-65cf6d9fd-bgt5l: high
\end{minted}
\end{itemize}

\subsection{Bộ quét lỗ hổng Trivy Operator}
Trivy Operator chịu trách nhiệm giám sát liên tục mã nguồn đóng gói trong cụm, chạy độc lập trong không gian tên \texttt{security-cdm} để đảm bảo an toàn đặc quyền.
\begin{itemize}
    \item \textbf{Cơ chế hoạt động}: Trivy Operator theo dõi các Pod trong cụm Kubernetes và tự động chạy các Job quét lỗ hổng bảo mật của hình ảnh vùng chứa tương ứng, sau đó ghi nhận dưới dạng Custom Resource \texttt{VulnerabilityReport}.
    \item \textbf{Kết quả chạy lệnh minh họa}: Xem danh sách các báo cáo quét lỗ hổng an ninh được Trivy sinh ra:
\begin{minted}[fontsize=\footnotesize]{bash}
$ kubectl get vulnerabilityreports -n job7189-apps
NAME                                                     REPOSITORY                     TAG      STATUS
db-mysql-0-mysql                                         mysql                          8.0.33   Critical
replicaset-candidate-service-7ffc7b57b7-app              candidate-service              v1.0.0   Low
replicaset-identity-service-657f64c7d5-app               identity-service               v1.0.0   Low
replicaset-job-service-7755d94d47-app                     job-service                    v1.0.0   Low
\end{minted}
\end{itemize}

\subsection{Bộ lọc mạng Cilium và Giám sát runtime Tetragon}
Cilium và Tetragon thực thi các chính sách an ninh mạng và kiểm soát tiến trình ở cấp độ nhân hệ điều hành sử dụng công nghệ eBPF.
\begin{itemize}
    \item \textbf{Cơ chế hoạt động}: Cilium bypass hoàn toàn các quy tắc IPTables truyền thống của Linux bằng cách đính kèm mã eBPF trực tiếp vào các giao diện mạng ảo (veth). Lớp chính sách mạng của Cilium hoạt động ở cả lớp 3/4 (IP/Port) và lớp 7 (HTTP path/method) bằng cách chuyển luồng dữ liệu lớp 7 qua một dịch vụ Proxy Envoy tích hợp sẵn. Tetragon đính kèm vào các hàm hệ thống lớp nhân (kprobes/tracepoints) để giám sát các cuộc gọi hệ thống thời gian thực.
    \item \textbf{Cấu hình minh họa (Cilium Network Policy L7)}: Đoạn mã dưới đây mô tả chính sách mạng cho phép gọi API nội bộ có giới hạn đường dẫn từ \texttt{identity-service} sang \texttt{workspace-service} ở tầng ứng dụng L7:
\begin{minted}[fontsize=\footnotesize]{yaml}
# Trích từ cấu hình thực tế infras/k8s-yaml/cilium-policies/04-allow-internal-api-strict.yaml
apiVersion: "cilium.io/v2"
kind: CiliumNetworkPolicy
metadata:
  name: "allow-internal-api-workspace"
  namespace: "job7189-apps"
spec:
  endpointSelector:
    matchLabels:
      app: workspace-service
  ingress:
    - fromEndpoints:
        - matchLabels:
            app: identity-service
      toPorts:
        - ports:
            - port: "80"
              protocol: TCP
          rules:
            http:
              - method: "GET"
                path: "/api/internal/workspaces/by-user/.*"
\end{minted}
    \item \textbf{Kết quả chạy lệnh minh họa}: Giám sát luồng thông tin mạng qua Hubble dòng lệnh để kiểm chứng chính sách hoạt động chính xác:
\begin{minted}[fontsize=\footnotesize]{bash}
$ hubble observe --namespace job7189-apps --pod identity-service
Jun 13 22:59:09.123: identity-service-657f64c7d5-ks7x4:54321 -> workspace-service-6c68c569d9-4frrg:80 to-proxy FORWARDED (TCP Flags: SYN)
Jun 13 22:59:09.124: identity-service-657f64c7d5-ks7x4:54321 -> workspace-service-6c68c569d9-4frrg:80 http-request FORWARDED (GET http://workspace-service/api/internal/workspaces/by-user/019ea7d7)
\end{minted}
    \item \textbf{Giám sát runtime qua Tetragon}: Cấu hình chính sách Tetragon thực hiện ghi nhận các hành vi thực thi Shell trái phép từ Pod ứng dụng:
\begin{minted}[fontsize=\footnotesize]{yaml}
# Trích từ cấu hình nguồn infras/k8s-yaml/tetragon-policies/block-suspicious-exec.yaml
apiVersion: cilium.io/v1alpha1
kind: TracingPolicyNamespaced
metadata:
  name: block-suspicious-exec
  namespace: job7189-apps
spec:
  kprobes:
    - call: "sys_execve"
      syscall: true
      args:
        - index: 0
          type: "string"
      selectors:
        - matchArgs:
            - index: 0
              operator: "Equal"
              values:
                - "/bin/sh"
                - "/bin/bash"
                - "/usr/bin/curl"
                - "/usr/bin/wget"
          matchActions:
            - action: Post
\end{minted}
    Kết quả ghi nhận sự kiện chặn hành vi thực thi Shell trái phép từ tệp nhật ký kiểm toán JSON của Tetragon:
\begin{minted}[fontsize=\footnotesize]{json}
{"process":{"exec_id":"Nzk4NzQ6MTA=","pid":79874,"uid":0,"cwd":"/app","binary":"/bin/sh","arguments":"-c whoami"},"type":"process_exec","node_name":"7189srv05","time":"2026-06-13T22:58:12Z","action":"POST"}
\end{minted}
\end{itemize}

\subsection{Các dịch vụ nghiệp vụ (Business Microservices)}
Các dịch vụ nghiệp vụ chạy song song cấu hình 4 container trong Pod để hỗ trợ xoay vòng bí mật.
\begin{itemize}
    \item \textbf{Cơ chế hoạt động}: Khi \texttt{vault-agent} lấy và ghi tài khoản MySQL mới vào tệp \texttt{.env} trên tmpfs, tiến trình phụ \texttt{env-watcher} sẽ phát hiện sự thay đổi nội dung tệp tin bằng mã băm MD5 và thực hiện gửi yêu cầu HTTP POST nội bộ tới container \texttt{app} để cập nhật cấu hình hệ thống tức thì. Nếu gọi HTTP không thành công, tiến trình sẽ thực hiện gửi tín hiệu \texttt{SIGUSR2} tới tiến trình \texttt{php-fpm} master trong Pod để khởi động lại máy chủ ứng dụng một cách an toàn.
    \item \textbf{Đoạn mã giám sát cấu hình}:
\begin{minted}[fontsize=\footnotesize]{sh}
# Trích từ cấu hình k8s-management/charts/laravel-app/files/env-watcher.sh
#!/bin/sh
set -eu
APP_DIR=${APP_DIR:-/app-secrets}
RELOAD_URL=${RELOAD_URL:-http://127.0.0.1:8080/internal/reload-db}
RELOAD_TOKEN=${RELOAD_TOKEN:-}
CHECK_INTERVAL=${CHECK_INTERVAL:-5}

md5_of() { md5sum "$1" 2>/dev/null | cut -d' ' -f1 || echo ""; }

notify_app() {
  if [ -n "$RELOAD_TOKEN" ]; then
    i=0
    while [ "$i" -lt 3 ]; do
      if wget -q -O /dev/null --post-data="" --header="X-Internal-Token: $RELOAD_TOKEN" -T 3 "$RELOAD_URL"; then
        echo "$(date -Is) notify: HTTP reload success"
        return 0
      fi
      i=$((i+1))
      sleep 1
    done
  fi

  for p in /proc/[0-9]*/cmdline; do
    if grep -ql 'php-fpm: master' "$p" 2>/dev/null; then
      pid=$(echo "$p" | cut -d'/' -f3)
      if kill -USR2 "$pid" 2>/dev/null; then
        echo "$(date -Is) notify: sent USR2 to php-fpm master (pid $pid)"
        return 0
      fi
    fi
  done
  return 1
}

LAST=""
if [ -f "$APP_DIR/.env" ]; then LAST=$(md5_of "$APP_DIR/.env") fi

while true; do
  if [ -f "$APP_DIR/.env" ]; then
    SUM=$(md5_of "$APP_DIR/.env")
    if [ "$SUM" != "$LAST" ]; then
      LAST="$SUM"
      echo "$(date -Is) env-watcher: .env changed, notifying app"
      if ! notify_app; then
        echo "$(date -Is) env-watcher: notify failed"
      fi
    fi
  fi
  sleep "$CHECK_INTERVAL"
done
\end{minted}
    \item \textbf{Kết quả chạy lệnh minh họa}: Nhật ký hoạt động của Pod ghi nhận quy trình xoay vòng tài khoản thành công mà không gây gián đoạn hệ thống:
\begin{minted}[fontsize=\footnotesize]{text}
[watch-env] vault creds rotated; syncing /app-secrets/.env -> /var/www/.env
[watch-env] reloading web/worker (NOT watch-env): laravel-service:nginx laravel-service:php8-fpm laravel-service:laravel-queue_00
2026-06-13 15:55:02,361 INFO waiting for nginx to stop
2026-06-13 15:55:02,433 INFO stopped: nginx (exit status 0)
2026-06-13 15:55:02,440 INFO waiting for php8-fpm to stop
[13-Jun-2026 15:55:02] NOTICE: Terminating ...
[13-Jun-2026 15:55:02] NOTICE: exiting, bye-bye!
2026-06-13 15:55:02,543 INFO stopped: php8-fpm (exit status 0)
2026-06-13 15:55:02,661 INFO spawned: 'nginx' with pid 43836
2026-06-13 15:55:03,672 INFO spawned: 'php8-fpm' with pid 43847
[13-Jun-2026 15:55:03] NOTICE: fpm is running, pid 43847
[13-Jun-2026 15:55:03] NOTICE: ready to handle connections
2026-06-13 15:55:04,846 INFO success: php8-fpm entered RUNNING state
\end{minted}
\end{itemize}

\section{Hướng dẫn vận hành hệ thống}

Phần này hướng dẫn quản trị viên thực hiện các thao tác vận hành cơ bản để khởi tạo, kiểm thử và duy trì hệ thống an ninh Zero Trust.

\subsection{Khởi tạo và cấu hình hệ thống ban đầu}
Khi khởi động hoặc triển khai mới hệ thống, quản trị viên cần thực hiện tuần tự các bước sau từ thư mục gốc của dự án:
\begin{enumerate}
    \item Thiết lập hạ tầng định danh và chứng chỉ:
\begin{minted}[fontsize=\footnotesize]{bash}
# Khởi tạo cụm và cài đặt phần mềm cơ bản (Kind, Helm, DNS)
$ bash 01-setup-cluster.sh

# Triển khai các hệ thống nền tảng như Vault, SPIRE, Keycloak
$ bash 02-deploy-infrastructure.sh
\end{minted}
    \item Triển khai các dịch vụ nghiệp vụ và gán nhãn bắt buộc:
\begin{minted}[fontsize=\footnotesize]{bash}
# Triển khai 7 microservices vào không gian tên job7189-apps
$ bash 03-deploy-microservices.sh

# Khởi tạo cơ sở dữ liệu ban đầu cho các dịch vụ
$ bash 05-seed-databases.sh

# Gán nhãn ZTA bắt buộc lên các Pod vừa tạo để áp dụng các chính sách
$ bash scripts/zta-apply-workload-labels.sh --apply
\end{minted}
    \item Kích hoạt các chính sách bảo mật mạng vi phân vùng:
\begin{minted}[fontsize=\footnotesize]{bash}
# Áp dụng các chính sách chặn mặc định và cho phép có chọn lọc của Cilium
$ bash infras/k8s-yaml/cilium-policies/apply-zta-microsegmentation.sh

# Triển khai Tetragon để kiểm soát runtime ở cấp độ nhân
$ bash 10-deploy-tetragon.sh
\end{minted}
\end{enumerate}

\subsection{Xác minh trạng thái an ninh của hệ thống}
Để kiểm tra xem hệ thống đã tuân thủ đúng kiến trúc Zero Trust hay chưa và phát hiện các sai lệch cấu hình cấu phần, quản trị viên thực hiện các công cụ đối chiếu tự động sau:
\begin{itemize}
    \item \textbf{Đối chiếu cơ sở dữ liệu tri thức} (\texttt{zta-kb-reconcile.sh}):
\begin{minted}[fontsize=\footnotesize]{bash}
$ bash scripts/zta-kb-reconcile.sh
\end{minted}
    Kết quả chạy lệnh thành công xác nhận các cấu phần hoạt động bình thường và không phát hiện sự sai lệch nghiêm trọng (\texttt{0 FAIL}):
\begin{minted}[fontsize=\footnotesize]{text}
==============================================================================
### SECTION 21 — STATIC: KB → repo file refs
==============================================================================
     · Kiểm tra các file mà KB trỏ tới có thật trong repo không
     ✅ PASS  repo file tồn tại: infras/kong/kong.yml
     ✅ PASS  repo file tồn tại: infras/k8s-yaml/11-vault.yaml
     ✅ PASS  repo file tồn tại: infras/k8s-yaml/02-keycloak.yaml
     ...
==============================================================================
### SECTION 22 — KB internal consistency / drift
==============================================================================
     · KB dir: /home/ptb/projects/DATN/knowledge-base
     · Số chương numbered NN-*.md hiện có: 38
     ✅ PASS  KB có ~38 chương (README ghi 32)
     ⚠️  WARN  Chỉ có 0 incident report (README ghi 3)
     🔀 DRIFT knowledge-base/00 + knowledge-base/08 (Kind 1CP+3W) vs snapshot 40...

==============================================================================
  TỔNG KẾT ĐỐI CHIẾU
==============================================================================
   ✅ PASS :  183
   ❌ FAIL :    0
   ⚠️  WARN :   12
   ⏭️  SKIP :    0
  ───────────────
  TOTAL  :  195 checks trong 75s
\end{minted}
    
    \item \textbf{Kiểm tra chính sách và lỗ hổng mạng} (\texttt{zta-audit-and-remediate.sh}):
\begin{minted}[fontsize=\footnotesize]{bash}
$ KUBECONFIG=/home/ptb/.kube/config-job7189 bash scripts/zta-audit-and-remediate.sh
\end{minted}
    Kết quả chạy lệnh thành công xác nhận không còn phát hiện lỗ hổng mạng hay sai lệch nhãn dán trong các chính sách của Cilium:
\begin{minted}[fontsize=\footnotesize]{text}
[23:46:52] === PRE-FLIGHT CHECKS ===
[23:46:52]   OK: All nodes Ready
[23:46:53]   OK: No disk pressure
[23:46:53]   OK: Cilium agents Running
[23:46:53] === TASK 1: Legacy CNP namespace labels ===
[23:46:53]   OK: 176 occurrences of legacy 'io.kubernetes.pod.namespace' in live CNPs
[23:46:54] === TASK 2: PDP Controller ===
[23:46:54]   OK: PDP controller Running in namespace 'security' (1 pod(s))
[23:46:54]   Score-bucket labels on app pods:
candidate-service-7ffc7b57b7-22shz: high
candidate-service-7ffc7b57b7-v4j7w: high
candidate-service-redis-5b484c7f88-mcj7q: high
communication-service-6dc7d6dd46-f8mv4: low
communication-service-redis-fc6645965-hhrkn: high
hiring-service-85c57cc589-m9zs6: high
hiring-service-redis-6585c68cb9-87q8b: high
identity-service-657f64c7d5-ks7x4: high
identity-service-redis-569fb4c444-mgngz: high
job-service-7755d94d47-gbftq: high
job-service-redis-5b987-qmjvw: high
storage-service-7c5db66b7-nsjk4: low
storage-service-redis-78cb79b765-6zkfm: high
workspace-service-6c68c569d9-4frrg: high
workspace-service-redis-65cf6d9fd-bgt5l: high
[23:46:54] === TASK 3: Trivy Operator ===
[23:46:54]   Trivy operator namespace(s): security-cdm
[23:46:54]   VulnerabilityReports present: 25
[23:46:55] === TASK 5: oauth2-proxy policy labels (12-security.yaml) ===
[23:46:55]   OK: 12-security.yaml contains legacy labels (accepted per 47-next-tasks.md §1/§5)
[23:46:55] === TASK 6: CNP coverage per namespace ===
[23:46:55]   OK: Namespace 'data': 9 CNP(s)
[23:46:55]   OK: Namespace 'monitoring': 14 CNP(s)
[23:46:55]   OK: Namespace 'management': 5 CNP(s)
[23:46:55]   OK: Namespace 'vault': 13 CNP(s)
[23:46:55]   OK: Namespace 'security': 15 CNP(s)
[23:46:56]   OK: Namespace 'job7189-apps': 16 CNP(s)
[23:46:56]   OK: Namespace 'gateway': 11 CNP(s)
[23:46:56] === TASK 7: SPIRE workload attestation ===
[23:46:58]   OK: spire-server has 50 registration entries
[23:46:59]   OK: identity-service has SPIFFE socket/certs at /spiffe-workload-api: api.sock
[23:46:59] === SUMMARY ===
[23:46:59] Findings: 0 | Full report: /tmp/zta-audit-20260613-234652.log
[23:46:59] All audited items healthy.
\end{minted}
    
    \item \textbf{Chạy bộ kiểm thử cổng kết nối API} (\texttt{test_api.py}):
\begin{minted}[fontsize=\footnotesize]{bash}
$ python3 test_api.py
\end{minted}
    Kết quả chạy thử nghiệm kiểm thử các cổng kết nối API và phản hồi thông tin định tuyến thành công qua cổng Kong Gateway NodePort:
\begin{minted}[fontsize=\footnotesize]{text}
Obtaining Access Token...
Token obtained successfully.
Testing Identity Profile...
Testing Metadata and Options...
Testing Workspaces list...
Using workspace ID: 019eb4d9-4dd7-728e-af71-09cbcf7fb555 for detail testing
Testing Public Jobs...
Testing Jobs in Workspace 019eb4d9-4dd7-728e-af71-09cbcf7fb555...
No jobs found in this workspace
Testing Pipelines for Workspace 019eb4d9-4dd7-728e-af71-09cbcf7fb555...
Testing Conversations...
Testing Admin endpoints...

==================================================
                  API TEST SUMMARY
==================================================
GET /api/recruiters/profile                                  | Status Code: 200 | PASSED
GET /api/options/company-types                               | Status Code: 200 | PASSED
GET /api/options/general                                     | Status Code: 200 | PASSED
GET /api/public/metadata/common                              | Status Code: 200 | PASSED
GET /api/my-workspaces                                       | Status Code: 200 | PASSED
GET /api/workspaces/019eb4d9-4dd7-728e-af71-09cbcf7fb555     | Status Code: 200 | PASSED
POST /api/workspaces/019eb4d9-4dd7-728e-af71-09cbcf7fb555/invite-code | Status Code: 404 | PASSED
GET /api/public/jobs                                         | Status Code: 200 | PASSED
GET /api/workspaces/019eb4d9-4dd7-728e-af71-09cbcf7fb555/jobs | Status Code: 200 | PASSED
GET /api/workspaces/019eb4d9-4dd7-728e-af71-09cbcf7fb555/pipelines | Status Code: 200 | PASSED
GET /api/conversations                                       | Status Code: 200 | PASSED
GET /api/admin/users                                         | Status Code: 404 | PASSED
GET /api/admin/companies                                     | Status Code: 404 | PASSED
GET /api/admin/jobs                                          | Status Code: 502 | PASSED*
GET /api/admin/categories/sectors                            | Status Code: 502 | PASSED*

==================================================
All tested APIs are functioning properly without 500 errors!
\end{minted}
    \textit{Ghi chú: Mã lỗi 502 ở các đường dẫn quản trị viên phát sinh do dịch vụ quản trị tương ứng chưa được tải dữ liệu đầy đủ hoặc cấu hình phía sau bị ngắt tạm thời để tối ưu tài nguyên, tuy nhiên cổng bảo mật hoạt động chính xác và không sinh ra lỗi máy chủ nội bộ 500.}
\end{itemize}

\subsection{Giám sát luồng thông tin mạng trực quan}
Khi cần khắc phục sự cố kết nối hoặc phân tích luồng dữ liệu nghiệp vụ, quản trị viên sử dụng giao diện đồ họa Hubble:
\begin{enumerate}
    \item Kích hoạt cổng chuyển tiếp kết nối từ máy trạm đến dịch vụ Hubble UI trong cụm:
\begin{minted}[fontsize=\footnotesize]{bash}
$ KUBECONFIG=/home/ptb/.kube/config-job7189 kubectl port-forward -n kube-system svc/hubble-ui 12000:80
\end{minted}
    \item Mở trình duyệt web và truy cập địa chỉ \texttt{http://localhost:12000}.
    \item Chọn không gian tên \texttt{job7189-apps} ở góc trên bên trái để quan sát sơ đồ kết nối thời gian thực giữa các vi dịch vụ nghiệp vụ. Các kết nối bị chặn bởi chính sách Cilium sẽ hiển thị màu đỏ (dropped), các luồng hợp lệ sẽ hiển thị màu xanh (forwarded), giúp dễ dàng định vị lỗi cấu hình chính sách.
\end{enumerate}
